%! Displaying and Describing Distributions — Unit 2, Lesson 2.1 — Homework
% This is the lesson's SCORED individual practice and the LAST component in the packet. Unit 2 on
% (2026-09-29): ASSIGNED at Close & Assign and done OUTSIDE class — not started in class. Due at
% the start of the first class after two study halls. Packet pages are the default; a Desmos
% activity may replace them on a given day (decided at Close & Assign).
% THIRD CONTEXT: the notes teach on Fox Run High, the Guided Practice runs on Pinecrest
% basketball, and these items are on the Oak Hollow Library summer reading challenge.
% TWO PAGES MAX: two boxes — items 1-2 on page 1; the histogram, items 3-6, and the spiralbox on
% page 2 — so no item breaks across a page.
% ITEMS: 1 core procedure (finish a stemplot whose frame is given, then read it); 2 the SPIRAL
% item (Unit 1: categorical vs quantitative — why genre cannot go in a histogram); 3 + 4 the
% CONTRAST PAIR (same questions; one distribution piles up low and is skewed RIGHT, the other
% piles up high and is skewed LEFT); 5 the CRUX head-on (Ava names the skew by the peak, the
% justify item); 6 interpret in context (full description with units).
% THE DATA — Oak Hollow Library, 20 kids.
%   Percent of reading goal met (sorted): 45 62 68 74 75 79 83 86 87 88 88 90 92 94 95 96 98 100
%     100 100. Stemplot missing leaves: 79, 88 (the second), 98. At least 90%: 9 kids.
%     Skewed left; peak 90–100; tail toward 45 (gap in the 50s). Median (88+88)/2 = 88.
%   Books read (histogram, bins of width 3): 0–2: 4, 3–5: 7, 6–8: 5, 9–11: 2, 12–14: 1,
%     15–17: 0, 18–20: 1 = 20. Skewed right; median in 3–5; gap 15–17; possible outlier 18–20.
% PAGE LOCKSTEP: the key differs only by saar-key, \blank -> \ans, \par\writeline ->
% \par\ansline (one for one — measured to land at the same height).
\documentclass[12pt]{article}
\usepackage{saar-article}
\usepackage{saar-boxes}
\newcommand{\TallMath}[1]{$\displaystyle #1\rule[-1.4em]{0pt}{3.2em}$}
% Word bank strip for every fill-in cluster. Defined per document; the key inherits it and
% prints the same strip, so heights match.
\newcommand{\wordbank}[1]{\par\vspace{2pt}\noindent\fcolorbox{goldacc}{goldbg}{\parbox{\dimexpr\linewidth-2\fboxsep-2\fboxrule\relax}{\small\textbf{Word bank:} #1}}\par\vspace{4pt}}

\begin{document}

\pageheader{Unit 2, Lesson 2.1}{Homework}
% NAMESTRIP: no \namedateperiod here. The student writes their name once, on the
% cover this component is stapled behind.

\vspace{0.12in}

% Authored identically in homework/ and homework_key/.
\begin{remindbox}
\small
\textbf{This is your graded homework.} It is scored out of 12 (2 points per item) and due at the
start of the first class after two study halls. It is assigned at the end of class and done on
your own --- the \emph{Keep in Mind} box on your cover is the page to flip back to.
\end{remindbox}

\vspace{0.08in}

\boxguard[20]
\begin{scenariobox}[Oak Hollow Library Summer Reading]{cerulean}
The Oak Hollow Public Library ran a summer reading challenge for $20$ kids. Each kid set a
reading goal. For every kid, the librarian recorded the \textbf{percent of the goal met} and the
\textbf{number of books read}.

\vspace{2pt}
Percent of goal met, in the order recorded: $88$, $100$, $74$, $95$, $62$, $90$, $83$, $100$,
$45$, $96$, $79$, $87$, $92$, $68$, $100$, $86$, $98$, $75$, $94$, $88$.
\end{scenariobox}

\vspace{0.08in}

\boxguard
\begin{notesbox}{Reading and Describing the Displays}
Every answer names the kids, the variable, and the units.
\wordbank{8 \;\textbullet\; 9 \;\textbullet\; 9 kids \;\textbullet\; quantitative \;\textbullet\; categorical}
\begin{enumerate}[leftmargin=1.6em, itemsep=9pt, topsep=3pt]

  \item The librarian's stemplot of percent of goal met is missing three values from the list.
        Fill in the three missing leaves, then read the display.

        \begin{center}
        {\large\renewcommand{\arraystretch}{1.3}%
        \begin{tabular}{r|l}
        4  & 5 \\
        5  & \\
        6  & 2\enspace 8 \\
        7  & 4\enspace 5\enspace \blank{0.7cm} \\
        8  & 3\enspace 6\enspace 7\enspace 8\enspace \blank{0.7cm} \\
        9  & 0\enspace 2\enspace 4\enspace 5\enspace 6\enspace \blank{0.7cm} \\
        10 & 0\enspace 0\enspace 0 \\
        \end{tabular}}
        \qquad
        \fbox{\textbf{Key:} $7 \mid 4 = 74\%$ of the goal}
        \end{center}

        How many kids met at least $90\%$ of their goal? \blank{2cm}

  \item The librarian also recorded each kid's favorite genre (mystery, fantasy, sports, and so on).
        Is favorite genre quantitative or categorical? \blank{3.4cm}

        \vspace{4pt}
        Could favorite genre be shown in a histogram? Explain.
        \par\writeline

\end{enumerate}
\end{notesbox}

\boxguard[30]
\begin{notesbox}{Reading and Describing the Displays, continued}
\begin{center}
\begin{tikzpicture}
  \node[font=\small\bfseries] at (3.85,3.2) {Books read this summer --- calculator histogram, bin width $3$};
  \foreach \y in {1,...,7} {
    \draw[linegray, very thin] (0,{\y*0.38}) -- (7.7,{\y*0.38});
    \node[left, font=\scriptsize] at (0,{\y*0.38}) {\y};
  }
  \foreach \x/\n in {0/4, 1/7, 2/5, 3/2, 4/1, 5/0, 6/1} {
    \fill[frostmid] ({\x*1.1},0) rectangle ({\x*1.1+1.1},{\n*0.38});
    \draw[cerulean, thick] ({\x*1.1},0) rectangle ({\x*1.1+1.1},{\n*0.38});
  }
  \draw[thick] (0,0) -- (7.8,0);
  \draw[thick] (0,0) -- (0,2.85);
  \foreach \x in {0,...,7} {
    \pgfmathtruncatemacro{\lab}{3*\x}
    \node[below, font=\scriptsize] at ({\x*1.1},0) {\lab};
  }
  \node[font=\scriptsize] at (3.85,-0.6) {books read};
  \node[rotate=90, font=\scriptsize] at (-0.6,1.4) {frequency};
\end{tikzpicture}
\end{center}
\wordbank{skewed left \;\textbullet\; skewed right \;\textbullet\; 3 to 5 books \;\textbullet\; 18 to 20 books \;\textbullet\; 90\% to 100\% \;\textbullet\; about 45\%}
\begin{enumerate}[leftmargin=1.6em, itemsep=7pt, topsep=2pt, start=3]

  \item \textbf{Books read} (the histogram). Where are most kids, and where does the long tail stretch?

        \vspace{4pt}
        Most kids at: \blank{3.6cm} \quad Tail toward: \blank{3.6cm}
        \par\vspace{6pt}
        Shape: \blank{3.6cm}

  \item \textbf{Percent of goal met} (your stemplot from item~1 --- its larger values are at the
        bottom). Where are most kids, and where does the long tail stretch?

        \vspace{4pt}
        Most kids at: \blank{3.6cm} \quad Tail toward: \blank{3.6cm}
        \par\vspace{6pt}
        Shape: \blank{3.6cm}

  \item Ava says the goal percents are skewed right, ``because most kids are up at the high end.''
        What did Ava name the skew by? Explain the correct shape using the tail.
        \par\writeline
        \par\writeline

  \item Describe the distribution of \textbf{books read}: shape, unusual features, center, and
        spread --- in context, with units.
        \par\writeline
        \par\writeline
        \par\writeline

\end{enumerate}
\end{notesbox}

\boxguard[5]   % the two-line spiralbox needs about five baselines; a larger guard pushes it to a third page
\begin{spiralbox}
\textbf{Next lesson: Measuring Center and Spread.} Today center and spread were ``about.'' Next
class the mean, the median, and the range make them exact --- and the shape you named today
decides which center to trust.
\end{spiralbox}

\end{document}
